\section{Typological transparency}\label{sec:6}

This section is methodological. The previous sections worked through cases in English; §\ref{sec:3} marks form--meaning applications without using them as evidence for a cross-domain profile. The question here is how to move from language-particular cases to a cross-linguistic claim. Haspelmath's \autocite*{haspelmath-2010-comparative-concepts} descriptive-versus-comparative-concept distinction is the right discipline.

\subsection{Haspelmath's distinction}

Haspelmath \autocite*{haspelmath-2010-comparative-concepts} argues that crosslinguistic comparison shouldn't assume universally available descriptive categories\footnote{Throughout, \mention{category} is used in Haspelmath's \mention{descriptive category} sense: a language-particular grammatical grouping fixed by a language's own distributional criteria, and neutral on whether that grouping is projectible. This differs from the sense used elsewhere in this research programme, on which a \mention{category} just is a projectible kind \autocite{reynolds2025hpcbook}; that kind-level notion is carried here by the \mention{projectibility profile} (§\ref{sec:7-1}), not by \mention{category}.} like ADJECTIVE, PASSIVE, ACCUSATIVE, or SUBJECT. Each language has its own descriptive categories, fixed by the language-particular criteria that make them work in that language \autocite[664]{haspelmath-2010-comparative-concepts}. Comparative concepts, by contrast, are analyst-created tools for comparison; they aren't part of any language system and can only be more or less well suited to the comparative task, rather than right or wrong \autocite[665]{haspelmath-2010-comparative-concepts}. They have to be built from universally applicable building blocks: conceptual-semantic notions, general formal notions, and other comparative concepts.

One pair makes the relation concrete. Haspelmath's comparative concept \enquote{dative case} is defined semantically, as a marker among whose functions is coding the recipient of a physical-transfer verb \autocite[666]{haspelmath-2010-comparative-concepts}. Language-particular categories then match it or not independently of their traditional names: the Finnish Allative matches, while the Nivkh Dative-Accusative doesn't, despite the name, because it codes a causee rather than a recipient \autocite[665--666]{haspelmath-2010-comparative-concepts}. Categorial particularism, Haspelmath's preferred position, holds that descriptive categories are language-particular and that comparison runs through analyst-defined comparative concepts standing in many-to-many relations with them. The position it rejects, categorial universalism, assumes a substantial set of crosslinguistic descriptive categories from which languages select \autocite[663]{haspelmath-2010-comparative-concepts}; \textcite{newmeyer-2007-crosslinguistic-formal-categories} is its canonical statement, taken up in §\ref{sec:6-objection}.

The methodological consequence is straightforward. A typological generalization isn't a claim about a descriptive category instantiated in every language. It's a claim about how analyst-defined comparative concepts are realized in language-particular ways. That nominative case-marking is widespread doesn't entail that \enquote{nominative} picks out the same descriptive category in any two languages.

\subsection{The umbrella schema as a comparative concept}

The umbrella schema is a comparative concept, not a descriptive category. It's built from the materials Haspelmath \autocite*[665]{haspelmath-2010-comparative-concepts} admits: \mention{property} and \mention{relation} are conceptual-semantic primitives, \mention{form} is a general formal notion. Asking the schema's question of a language isn't claiming that the language has a category corresponding to \term{transparency}.

\mention{Accessible through} is the contested slot, since accessibility has discipline-specific senses (discourse-pragmatic salience, psycholinguistic retrievability, formal-semantic substitution-licensing). The schema's sense is more abstract: $P$ is accessible to $R$ through $X$ when $R$ tracks $P$ across $X$ rather than tracking a property of $X$ itself. The case studies operationalize that differently, by agreement targeting, constructional interpretability, rated transparency, or priming magnitude, but the comparative concept lives above the operationalizations.

The non-instance test (§\ref{sec:2-3}) is what keeps the concept honest, because its counterpart is fixed language-internally. The test asks whether the language under study attests a same-type variant of $X$ whose structure makes $P$ inaccessible to $R$. For English number-transparent NPs that counterpart is head-driven agreement, and the transparent class is the marked exception against it; the pattern isn't exported. The procedure is identical in any language: identify $X$, $P$, and $R$ by that language's own distributional criteria, then look for the concealing variant. Where it exists, the transparent case is informative; where the construction-type admits none, as with English coordination, the schema returns a non-instance. So \mention{accessible through} imports no English-particular grammar. It inherits whatever counterpart the local construction-type supplies.

So the bindings are language-particular while the question isn't. In English, $X$ is the quantificational noun plus \mention{of}-complement construction, a \textit{CGEL} descriptive category; $P$ is the complement's number/count profile; $R$ is subject-verb agreement. In the worked cases below, Spanish partitive subjects and Tsez complement clauses fill the slots by different language-particular criteria. That's what makes the typological question tractable: not \enquote{do all languages have number-transparent constructions?}, which would look for English-style NPs everywhere, but \enquote{which configurations in this language let a property remain accessible to a relation that the surface form might have hidden?} The answers are expected to differ, and the English, Spanish, and Tsez cases aren't instances of one descriptive category.

Two limits on what that buys. The claim at the comparative-concept level is structural, that the concept lives above the descriptive category whatever the language, rather than a demonstrated comparative result: the extensions on hand are the Spanish and Tsez cases below. The form--meaning family (§\ref{sec:3}) isn't itself a comparative concept; the schema is, and the family is a set of language-particular phenomena and measurements compared through it.

\subsection{Cross-linguistic cases and topologies}\label{sec:6-topologies}

Two worked languages let the question be asked outside English, one of them outside the agreement type English uses. Spanish has an analogous agreement-access pattern. In a partitive subject like \mention{la mayoría de los manifestantes} (`the majority of the demonstrators'), the finite verb may agree in the plural with the complement or in the singular with the quantificational noun \mention{mayoría}, as in (\ref{ex:spanish-mayoria}).

\ea\label{ex:spanish-mayoria}
\ea \gll La mayoría de los manifestantes gritaban. \\
the majority of the demonstrators shout.\textsc{ipfv.3pl} \\
\glt `The majority of the demonstrators were shouting.'
\ex \gll La mayoría de los manifestantes gritaba. \\
the majority of the demonstrators shout.\textsc{ipfv.3sg} \\
\glt `The majority of the demonstrators was shouting.'
\z
\z

The plural (\ref{ex:spanish-mayoria}a) is the more frequent option and the only one available under a plural predicative \autocite{rae-concordancia-mayoria}. Both variants are attested on the one subject, much as English \mention{majority} and \mention{minority}, \mention{mayoría}'s boundary cognates, take occasional singular agreement alongside the dominant plural override (§\ref{sec:2-2}). Spanish has its own literature on the alternation, which treats it as \term{concordancia ad sensum} and reaches two conclusions that matter here. First, both options are genuine agreement rather than agreement displaced by government, so the plural isn't a species of \term{silepsis} \autocite{sanjulian2018-cuantificadores}. Second, the alternation isn't free across syntactic contexts, and singular agreement with the quantifier noun is compatible with a distributive reading of the clause \autocite{sanjulian2018-cuantificadores}, which blocks the inference that notional plurality forces the override. The partitive and pseudo-partitive structures the alternation runs on are themselves heterogeneous in Spanish \autocite{sanjulian2018-pseudopartitivas}, as they are in English (§\ref{sec:2-1}).

What travels to Spanish is the access pattern, not the English cluster: no Spanish determiner, complementation, ellipsis, or premodification profile has been shown here. And the concealment control is weaker than the English one, since the singular option with \mention{mayoría} is dispreferred rather than blocked.

Tsez (Nakh-Daghestanian; head-final and ergative; about seven thousand speakers) gives a case in a different marking type. Tsez verbs prefix-agree in noun class with their absolutive argument, so the agreement is marked on the verbal head. A clausal argument is class IV, the class that includes clauses, so a matrix verb's default is class IV agreement with the complement clause as a unit; but when the embedded absolutive is a topic, the matrix verb can instead agree with it \autocite[436]{potsdam-polinsky-1999}.

\ea\label{ex:tsez-lda}
\ea \gll enir [už-ā magalu b-āc'ru-łi] r-iyxo \\
mother [boy-\textsc{erg} bread.\textsc{iii.abs} \textsc{iii}-ate-\textsc{nmlz}] \textsc{iv}-knows \\
\glt `The mother knows the boy ate the bread.' (default: the matrix verb is class IV, agreeing with the clause)
\ex \gll enir [už-ā magalu b-āc'ru-łi] b-iyxo \\
mother [boy-\textsc{erg} bread.\textsc{iii.abs} \textsc{iii}-ate-\textsc{nmlz}] \textsc{iii}-knows \\
\glt `The mother knows the boy ate the bread.' (long-distance agreement: the matrix verb is class III, agreeing with embedded \mention{magalu} `bread')
\z
\z

In the schema's terms, $X$ is the embedded clause, $P$ is the noun class of the embedded absolutive, and $R$ is matrix verb agreement. The concealing variant, and here genuinely the default, is the class IV agreement in (\ref{ex:tsez-lda}a), attested in the same language and the same sentence, and (\ref{ex:tsez-lda}b) is the transparent override. The counterpart isn't \enquote{the nearest head's own feature}, as with English number-transparent nouns, but \enquote{default class IV for a clausal argument}: a different default in a different language type. That's what the comparative concept predicts, the question constant and the language-particular answer not. Whether the override is analysed as covert movement or as long-distance agreement is a framework-internal question the schema brackets, as with extraction in §\ref{sec:2}; the fuller analysis is in \textcite{polinsky-potsdam-2001}.

Spanish and Tsez are both outward endpoint cases: a property associated with material inside $X$ reaches an external agreement relation. The same schema maps several other accessibility topologies, provided each remains subject to the §\ref{sec:2-3} counterpart test: identify $X$, $P$, and $R$, then ask what same-language pattern would make the proposed access disappear. Appendix~\ref{sec:B} sets out four of them with their applicability and status conditions (null-source subpart access in Aleut, inward dependent access in Lardil and Kayardild, path-distributed access in Chamorro, and the feature-relay boundary case in Kutenai and Algonquian obviation), along with a targeted check of head-marked systems.

Two results from that survey bear on the argument here. Aleut is a qualified positive: the controller of $R$ can be recoverable as a subpart of a non-subject NP without being overt, though Merchant's analysis treats the unpronounced possessor as syntactically present pro, which keeps the case distinct from overt intact-NP transparency \autocite[195]{merchant-2011-aleut-case-matters}. And the head-marked candidates cluster around possessors, most of them analysed as externalized, by possessor raising, applicative mediation, or external possession, rather than as agreement with a dependent that stays NP-internal \autocite{ritchie-2016-prominent-internal-possessors,deal-2017-external-possession}. So the diagnostic point is narrow: overt sensitivity to a possessor isn't enough, and the decisive question is whether an embedded dependent remains NP-internal while controlling clause-level morphology.


\subsection{Mapping transparency and access transparency}\label{sec:6-mapping}

There's an established typological programme under the name \term{transparency}, and it classifies the paper's paradigm cases as non-transparent. Setting the two side by side is the clearest available demonstration that \term{transparency} isn't one projectible kind.

Leufkens defines transparency as \enquote{a one-to-one relation between meaning and form}, treating every structure that violates such a relation as non-transparent \autocite[2]{leufkens2015}, and ranks 22 languages by how many non-transparent features their grammars show; the framework is extended to a larger sample and an implicational hierarchy in \textcite{hengeveld-leufkens-2018}. Opacity divides into four categories according to how the one-to-one relation fails \autocite[16, §§4.1--4.4]{leufkens2015}: redundancy, discontinuity, fusion, and form-based form. Agreement is the leading instance of the first. Where one unit of meaning is expressed by several forms, \enquote{one of the forms is not necessary; it is redundant, as it does not provide any additional information}, and clausal and phrasal agreement are catalogued in that category \autocite[18, §§4.1.3--4.1.4]{leufkens2015}.

That's the opposite verdict from this paper's on the same configurations. Subject--verb agreement is the paradigm $R$ for which a property remains accessible through a mediator; on the one-to-one conception it's redundant marking, and a language with more of it is less transparent. Both verdicts are correct within their own question. They're different questions: whether a language's levels map one-to-one, and whether a particular relation can reach a property across a mediator that would otherwise conceal it. Call the first \term{mapping transparency} and the second \term{access transparency}.

English number-transparent QNs cross-classify. They're an access-transparency instance. Being agreement, they're also a mapping-transparency violation. A language could reduce its agreement, becoming more transparent in Leufkens's sense, and thereby lose the very relations that make access transparency visible. The divergence isn't a disagreement to be resolved by defining terms more carefully, because the two programmes individuate by different relata: levels of organization in one, mediator-property-relation triples in the other. An umbrella that delivered opposite verdicts on one construction while both verdicts stood wouldn't be a kind, and that's what §\ref{sec:7} concludes on independent grounds.

Two further payoffs. Leufkens's inventory is organized, so it supplies a principled comparison class rather than another loose sense of the word, and the redundancy/discontinuity/fusion/form-based-form taxonomy is a checklist this paper's cases can be run against. And her finding that every language sampled has some non-transparent features, most often of the redundancy type \autocite{leufkens2015}, means agreement relations are common, so they're where a search for mediated access is most likely to find candidates. That says nothing about how often the search succeeds.

\subsection{An objection and a reply}\label{sec:6-objection}

A reader might object that the umbrella schema looks like a universal descriptive category smuggled in as a comparative concept. The sharpest form of the objection targets the non-instance test: \mention{accessible through} presupposes a concealing counterpart, and for the English cases that counterpart is head-driven agreement, which is independently the unmarked pattern, a language-particular fact. If the schema's deepest primitive is parasitic on an English-shaped default, then calling it universally applicable looks like categorial universalism in disguise.

The reply isn't that the schema is \enquote{just a question}: a universal category can be framed as a question too, so that move wouldn't tell a comparative concept apart from a category in disguise. The reply is that the counterpart the objection points to is fixed \mention{language-internally}, which is how comparative concepts are meant to work. The concealing counterpart is not assumed to be English head-driven agreement; it's whatever same-type variant the language under study attests, by the procedure set out earlier in this section. So the schema doesn't import an English default; it inherits each language's own.

Three things follow that a universal descriptive category wouldn't deliver. The variables are bound by language-particular distributional criteria, not by a fixed cross-linguistic definition. The schema stands in many-to-many relations with descriptive categories: one language's category can fill $X$, $P$, or $R$ under different bindings, and several categories can realize one binding. It also carries explicit status conditions: a candidate is a non-instance where the construction-type lacks the required mediator or bearer, and a profile claim is demoted where the predicted access fails to project, as in the \mention{bunch} localization of §\ref{sec:2-2}. The honest residue is that the test is only as sharp as the prior individuation of construction-type. Deciding what counts as the same type with and without concealment is itself a language-internal analysis, so the diagnostic is crisp where that individuation is independently motivated and softer where it isn't. A candidate counterpart that isn't independently motivated can't supply a confident positive. The procedure's refusal to overgeneralize is a useful stress test, though indeterminacy alone doesn't validate the test. Haspelmath's \autocite*[665]{haspelmath-2010-comparative-concepts} criterion for a good comparative concept, that it can be more or less well suited to the task of permitting crosslinguistic comparison, is exactly the standard the umbrella schema is meant to meet.

Newmeyer's specific argument deserves a direct response. His case is that typological generalizations must reference not just a universal concept but the specific form realizing it: negative particles and negative auxiliaries express the same concept, but only the latter pattern with verbs in word-order typology, so a generalization stated over \enquote{negation} misses the regularity \autocite[136]{newmeyer-2007-crosslinguistic-formal-categories}. The point is correct on its own terms. What follows from it is weaker than crosslinguistic \mention{descriptive} categories, because the formal reference it demands enters at the $X$ slot, and the slot and its values are different things. The slot is \mention{typed} with general-formal notions in Haspelmath's \autocite*[666]{haspelmath-2010-comparative-concepts} sense: construction, clause, NP, dependency domain. A value of the slot is a language-particular configuration fixed by that language's criteria, the number-transparent QN construction in English and the configurations the analyst identifies in each comparison language. The same holds for $P$ and $R$: the slots are abstractly typed, while their values are particular number categories, grammatical relations, and agreement systems. Nothing here requires equating configurations token-for-token across languages. The concealing counterpart that gives the non-instance test its grip is itself such a general-formal concept, not an imported universal category.

The metaphysical conclusion in §\ref{sec:7}, that the umbrella schema is a tool for locating projectibility profiles rather than itself a kind, is consistent with this. Three levels stay distinct: the comparative concept (mediated accessibility) lives at the cross-linguistic methodological level; descriptive categories (English number-transparent quantificational nouns, Spanish partitive subjects, Tsez long-distance agreement, and so on) live at the language-particular grammatical level; \term{projectibility profiles} are observable bases licensing non-trivial projection, stabilized by identifiable sources, realized within or across such descriptive categories. Kindhood is a downstream verdict on profiles that earn it. The paper doesn't collapse the levels.

Two limits on the position taken here. Haspelmath's distinction is itself contested and this paper doesn't adjudicate; it's adopted as the safer default, since it requires no commitment to descriptive-category universality, and a reader more comfortable with categorial universalism can take the schema as a working concept whose status is secondary to its methodological pay-off. And Croft's \autocite*{croft-2001-radical-construction-grammar} Radical Construction Grammar would route the same commitment through construction-similarity hierarchies rather than comparative concepts; the schema is meant to sit above any single construction-grammar variant, so that route is left open rather than argued against.

One consequence carries into §\ref{sec:7}: cross-linguistic projection at the profile level isn't entailed by cross-linguistic projection at the comparative-concept level. Asking the question everywhere doesn't make the answers travel.
